\chapter{捧腹大笑的机制分析：从五种解释到可检验假设}\label{ch:07}

\begin{nlkeypoints}
  \begin{enumerate}[label=(\arabic*),leftmargin=2.4em]
    \item 第一版教材的五重机制（情绪代偿、群体稀释、反精致、校园社交、
          逆反认证）被重写为\keyw{“压力源—补偿行为—可观察指标—
          竞争性解释”四列矩阵}，使其从感想变为可检验的分析框架
          （矩阵为本书构造，\conf{M} 级）。
    \item 五个机制各有学理来源（伯格林的机械性、弗洛伊德的能量节省、
          巴赫金的降格、普罗文的社交笑、
          Nissenbaum--Shifman 的文化资本），本章逐一给出
          \keyw{可检验假设与证伪条件}（H7.1--H7.5）。
    \item 机制之间不是并列关系而是\keyw{分层关系}：
          M1（代偿）与 M2（稀释）作用于个体—情绪层，
          M3（反精致）作用于风格层，M4（校园社交）作用于传播层，
          M5（逆反认证）作用于权威关系层；
          每层提出的指标不同，因此五个机制\keyw{可以被分别证伪}，
          而不是一荣俱荣。
  \end{enumerate}
\end{nlkeypoints}

本章处理全书的核心谜题（问题 Q2）：\keyw{为什么大学生会笑，
而且愿意反复转发}。第一章已把“为什么聪明人爱看无厘头”
改写为三种竞争读法，本章的任务是把其中功能—符号读法的
五个候选机制逐一学理化：先给每个机制找一个成熟的理论来源，
再把“机制”写成“压力源—补偿行为—可观察指标—竞争性解释”
的四列结构，最后为每个机制提出假设与证伪条件。
必须预先声明：本章没有任何经验数据，全部假设处于
\keyw{待检验}状态（第十一章给出检验工具）。
本章的分析对象——“捧腹大笑”这一表达形式——
其符号结构已在第六章给出；本章问的不是它是什么，
而是它\keyw{在做什么、为什么有人愿意反复做}。

\section{从“五重战神”到机制矩阵}\label{sec:07:matrix}

第一版教材第五部分以“捧腹大笑思想理论体系（迫真）”为题，
列出了五个机制并各给一段评述，末尾自嘲“（不行，扯不下去了）”。
这份材料的直觉质量很高——五个机制的\keyw{命名}在社群传播中
已经被反复引用——但它有三个方法论缺陷：

\begin{enumerate}[label=\arabic*.]
  \item \textbf{平行罗列}：五个机制之间是什么关系（互补？竞争？
        分层？），原文没有回答；
  \item \textbf{不可证伪}：每段评述都是“机制名 + 现象复述”，
        没有给出“如果该机制不成立，我们会看到什么”；
  \item \textbf{层次混淆}：情绪机制（代偿）与传播机制（校园社交）、
        权威关系机制（逆反认证）被放在同一层并列，
        导致“一个机制反例”会被误读为“整个体系崩塌”。
\end{enumerate}

本书的重写方案是把每个机制装进同一个四列结构：

\begin{table}[htbp]\centering
  \caption{五重机制的分析矩阵模板（本书构造的教学框架，\conf{M} 级）}
  \label{tab:07:matrixtpl}
  \small
  \begin{tabularx}{\textwidth}{@{}lX@{}}
    \toprule
    列 & 内容与要求 \\ \midrule
    压力源 & 机制所应答的现实处境；必须是可指认的情境
      （如“绩点竞争”），不得是“当代社会”这类不可操作的笼统表述 \\
    补偿行为 & 使用者实际做的事（转发、配文、立牌、摆架势）；
      必须与第六章立场维度的语用行为对应 \\
    可观察指标 & 单一数字或单一频率（如“压力考试周内的转发占比”）；
      必须原则上可计数 \\
    竞争性解释 & 若不引入该机制，哪些更简单的解释也能覆盖同一指标；
      机制只有在排除竞争解释后才有存在必要 \\ \bottomrule
  \end{tabularx}
\end{table}

四列的次序就是论证的次序：压力源立“需求”，
补偿行为立“供给”，可观察指标立“检验”，
竞争性解释立“对手”。任何一个机制若填不满四列
（尤其是填不出可观察指标），就不配称为机制，
只是修辞。后文各节按此模板逐个展开。

\section{笑的学理基础：四个古典来源}\label{sec:07:theory}

在进入机制之前，先把“笑”本身的理论资源铺开。
四个来源分别回答四个问题，后文机制将分别调用。

\subsection{优越论与机械论：笑作为社会矫正}

伯格森在《笑》中提出著名的公式：
\keyw{“镶嵌在生命之上的机械性”是可笑的}，
笑是“机械性对生命性的报复”——
当一个生命体表现得像机器（僵硬、重复、自动化）时，
旁观者的笑执行了一种社会矫正功能\cite{bergson1900rire}。
奶蛙语料中“僵直的爬行、重复的蹲起、程式化的摆架势”
（第六章衍生动作）正是“机械性”的具象化；
而 AI 生成的畸变身体（膨腹、细肢）则是
“生命形态的机械故障”的极端版本。
伯格森框架的预测：\keyw{笑点应随“机械化程度”单调增强}——
这是一个可操作化的变量（动作重复次数、程式化程度编码）。

\subsection{释放论：笑作为能量节省}

弗洛伊德在《诙谐及其与无意识的关系》中把笑解释为
\keyw{心理能量的节省}：诙谐绕过了压抑与批判性检查，
原本用于抑制的心理能量被省下、以笑的形式释放
\cite{freud1905witz}。该框架对奶蛙语料的对应物是
“绷不住”体验：紧绷（抑制）与绷不住（释放）的对立
恰好是能量模型的大众版。释放论的预测：
\keyw{禁忌或高压力语境中的笑应更强烈}——
第七章 M1（情绪代偿）直接调用这一预测。

\subsection{狂欢论：笑作为秩序倒置}

巴赫金把狂欢节式的笑描述为\keyw{降格}
（\en{degradation}）：一切崇高、精尖、正经之物
被拉向身体下部与物质底层，在笑声中完成暂时的等级翻转
\cite{bakhtin1968rabelais}。狂欢论的预测：
笑的对象应系统性地偏向\keyw{权威与精致符号}，
而非随机对象——第七章 M3（反精致）与
M5（逆反认证）分别调用这一预测的两半。

\subsection{社交论：笑作为关系信号}

普罗文的实验研究表明，笑远不只是“好笑”的情绪反应：
它是高度定型化的社交发声，绝大多数笑发生在社交情境中，
且与逗笑内容的质量弱相关；笑具有强传染性
\cite{provine1992laughter,provine2000laughter}。
社交论的预测：\keyw{同一素材的笑应随共同在场人数增加而增强}——
第七章 M4（校园社交）调用这一预测，
第十章的宿舍—教室场景为其检验场。

四个来源的分工可以压成一句话：
伯格森解释\keyw{笑什么}（机械性），
弗洛伊德解释\keyw{为什么释放}（能量节省），
巴赫金解释\keyw{冲着谁笑}（等级倒置），
普罗文解释\keyw{和谁一起笑}（社交信号）。
五个候选机制各自从中取用一角，
但都不与任何一个理论画等号——
理论提供的是\keyw{可检验预测的来源}，不是结论本身。

\begin{nlcritique}[批判性反思：理论的移植风险]
  四个古典理论各有其历史语境：伯格森谈的是沙龙与剧场的
  喜剧观众，弗洛伊德的材料是十九世纪末维也纳的诙谐，
  巴赫金的狂欢节是中世纪广场文化，普罗文的实验是
  北美大学生。把它们移植到“短视频平台的算法分发”
  语境，每一个都存在\keyw{前提错位}的可能——
  例如，算法推荐改变了“共同在场”的含义，
  普罗文的社交笑结论未必平移成立。
  因此本章的使用方式是\keyw{把理论当假设生成器}，
  而非当证据：每个理论只贡献一条可检验预测，
  预测不成立时被淘汰的是“理论在奶蛙语境的应用”，
  不必连理论本身一起否定。
\end{nlcritique}

\section{M1 情绪代偿：把“绷不住”合法化}\label{sec:07:m1}

\subsection{机制表述}

第一版教材的原文：“考研、绩点、就业的复合压力下，
负面情绪缺少体面的出口。奶蛙提供了一种‘笑着崩溃’的姿态
——既表达了失控，又不至于沉重。”
按四列矩阵重写：

\begin{table}[htbp]\centering
  \caption{机制 M1（情绪代偿）的四列展开（\conf{M} 级）}
  \label{tab:07:m1}
  \small
  \begin{tabularx}{\textwidth}{@{}lX@{}}
    \toprule
    压力源 & 绩点竞争、考研备考、求职季等可指认的高压情境
      （以学期周期为节律） \\
    补偿行为 & 在压力情境的社媒表达中使用奶蛙大笑素材，
      以“笑”的语体发出原本沉重的自述 \\
    可观察指标 & 压力事件（期中考、考研倒计时节点）前后，
      奶蛙类内容使用频率的\keyw{相对变化率}；
      与一般娱乐内容使用频率变化率之差 \\
    竞争性解释 & (a) 纯消遣：高压周里人们本来就刷更多手机，
      奶蛙只是“恰好刷到”；(b) 时间填塞：备考期社交时间减少，
      低成本内容占比被动上升 \\ \bottomrule
  \end{tabularx}
\end{table}

\subsection{学理对应与假设}

情绪代偿的学理骨架是弗洛伊德的释放论
（\ref{sec:07:theory} 小节）加上情绪的“行动倾向”理论：
Frijda 主张情绪内置行动准备（emotional action tendency），
表达行为是行动倾向的出口\cite{frijda1986emotions}。
“笑着崩溃”在这个框架里读作：
崩溃的行动倾向（退出、哭、求助）被“笑”这一
社交许可的出口替代，且替代不损失信息量——
接收方仍能读出崩溃（第六章立场维度的“自陈替代”）。

\begin{hypothesis}[H7.1 情绪代偿假设]\label{hyp:07:m1}
  在高校用户的社媒使用中，奶蛙类内容的使用频率
  随压力情境的出现而\keyw{选择性上升}：
  上升幅度显著高于其总娱乐内容使用频率的上升幅度，
  且该差值在自评压力更高的用户中更大。
\end{hypothesis}

\keyw{证伪条件}：若压力期与非压力期的奶蛙使用频率比，
不高于总娱乐内容的两期频率比（即无选择性），
或者选择性上升与自评压力无关，则 H7.1 被证伪，
情绪代偿退化为“时间填塞”的被动效应。
\keyw{必须声明}：H7.1 尚无任何经验检验，
第十一章的问卷设计（量表 + 学期节律对齐）是为其准备的。
“笑着崩溃”的表达门槛问题还与在线去抑制效应相关：
匿名与不可见性降低情绪表达成本
\cite{suler2004disinhibition}——这一因素在
H7.1 中被当作常量背景，不作单独检验。

\section{M2 群体稀释：用“我们”消解“我”}\label{sec:07:m2}

\subsection{机制表述}

第一版教材原文：“用一张梗图代替自述，把‘我很差劲’
翻译成‘我们都这样’。个体暴露被群体化收编，
自嘲因而不必付出羞耻成本。”重写：

\begin{table}[htbp]\centering
  \caption{机制 M2（群体稀释）的四列展开（\conf{M} 级）}
  \label{tab:07:m2}
  \small
  \begin{tabularx}{\textwidth}{@{}lX@{}}
    \toprule
    压力源 & 自我暴露的羞耻成本：直接陈述失败或焦虑
      需要承担社交评价风险 \\
    补偿行为 & 以梗图、群发、评论区复读等方式发布，
      使个人陈述获得“大家都这样”的群体外观 \\
    可观察指标 & 第一人称自述句与奶蛙梗图配对出现的比例；
      同一内容的“带梗图 / 不带梗图”两版获得的
      自我暴露类回复数之差 \\
    竞争性解释 & (a) 表达习惯：梗图只是流行文体，与羞耻无关；
      (b) 效率：图像比文字输入更快，用户选图仅为省事 \\
    \bottomrule
  \end{tabularx}
\end{table}

\subsection{学理对应与假设}

M2 的学理骨架是戈夫曼的拟剧论与情绪感染的合流。
戈夫曼的\keyw{前台/后台}区分指出，个体在前台
管理自己给他人留下的印象\cite{goffman1959presentation}；
梗图正是前台管理装置——情绪信息（后台事实）
经由一个“非我”的图像发出，发送者保持前台体面。
情绪感染研究则说明接收端如何“接住”这一信号：
情绪在群体中经由模仿与反馈循环传播
\cite{hatfield1994contagion}，数字环境放大了
发送者—媒介—接收者三要素中的媒介作用
\cite{goldenberg2020digital}。“我们都这样”之所以能成立，
正因为梗图天然携带\keyw{可复读性}——
它邀请接收者以转发代替安慰，
把一次情绪事件转写为一次集体复读。

\begin{hypothesis}[H7.2 群体稀释假设]\label{hyp:07:m2}
  同一负面自述内容，附带奶蛙类梗图时：
  (a) 发布者事后删除率更低（羞耻成本更小的行为学代理指标）；
  (b) 收到的回复中“同感/复读类”占比更高、“安慰/追问类”占比更低。
\end{hypothesis}

\keyw{证伪条件}：若带梗图与否不影响上述两项指标，
则“稀释”只是文体偏好，不承担心理功能。
M2 与第五章“文化资本”框架的连接点是：
用梗代替自述同时完成\keyw{情绪表达}与
\keyw{身份认证}两个动作\cite{nissenbaum2015contested}——
但必须指出，身份认证也能解释同样的行为
（用梗因为“圈内人应该用梗”），这正是 M2 的
竞争性解释之一，判别设计见 \ref{sec:07:compete} 节。

\section{M3 反精致：对优秀人设的审美罢工}\label{sec:07:m3}

\subsection{机制表述}

第一版教材原文：“大学生被期待为‘精致的、上进的、
有规划的人’。转发一只丑、怪、抽象的奶蛙，
是一次低成本的不合作表态。”重写：

\begin{table}[htbp]\centering
  \caption{机制 M3（反精致）的四列展开（\conf{M} 级）}
  \label{tab:07:m3}
  \small
  \begin{tabularx}{\textwidth}{@{}lX@{}}
    \toprule
    压力源 & “精致人设”的维持成本：朋友圈、社交主页
      中的上进形象要求与其实际状态的落差 \\
    补偿行为 & 主动转发畸变的、劣化的、反审美的形象，
      以展示对精致规范的不合作 \\
    可观察指标 & 用户内容流中“故意劣化”类内容占比；
      同一用户在成绩/求职等评价敏感期的劣化内容占比变化 \\
    竞争性解释 & (a) 审美多元：劣化风格本身即审美偏好，
      无表态含义；(b) 模仿流行：跟随社群风格而非对抗规范 \\
    \bottomrule
  \end{tabularx}
\end{table}

\subsection{学理对应与假设}

M3 直接调用巴赫金的降格概念
（\ref{sec:07:theory} 小节）：把“精致”拉入
“丑怪—劣化—荒诞”的物质层面，是狂欢式倒置的当代版。
伯格森的机械论也在此处工作：
奶蛙的“机械性身体”（重复动作、畸变比例）
与精英人设所要求的“得体的生命管理”形成对照，
笑它即是宣告不参与该规范\cite{bergson1900rire}。
亚文化研究为“不合作表态”提供了更精确的表述：
Hebdige 指出亚文化风格是通过\keyw{符号层面的抵抗}
组织起来的，而抵抗终将被商业与主流收编
\cite{hebdige1979subculture}——后者正是第九章
（奶蛙周边商业化）要展开的“收编”环节。

\begin{hypothesis}[H7.3 反精致假设]\label{hyp:07:m3}
  在评价敏感期（绩点公布、求职结果发布），
  用户发布“故意劣化”内容的频率相对基线的上升幅度，
  显著高于非敏感娱乐内容的同期上升幅度；
  且该上升在自我呈现压力更高（如以学术/精英形象为主页）的
  用户中更强。
\end{hypothesis}

\keyw{证伪条件}：劣化内容使用在敏感期与非敏感期无差异，
或上升与主页人设类型无关。M3 的一个独特困难是
“故意劣化”的编码判定——需要区分\keyw{真的做不好}
与\keyw{故意做坏}，编码手册必须给出
（如“刻意使用低分辨率素材、故意错位排版、
拒绝滤镜修正”等可观察标记），这一编码方案
在第十一章给出。

\begin{nlcritique}[批判性反思：“反精致”的阶级自嘲问题]
  “反精致”的表述有一个必须直面的张力：
  第一版教材的读者群（顶尖高校学生）恰好是
  “精致竞争”的\keyw{优胜者或准优胜者}。
  “优胜者扮演反规范”与传统亚文化中
  “弱势者表达抵抗”的结构并不相同——
  前者的“罢工”没有实际代价，甚至可能成为
  新一轮身份竞赛的筹码（“我反内卷”本身成为人设）。
  这提示 M3 可能只是\keyw{符号消费}而非\keyw{规范抵抗}：
  判别依据是“反精致表态是否伴随实际行为让步”
  （如真的退出绩点竞争）。若无伴随行为，
  M3 应降级为“人设的一部分”。
  本书不预设哪一种成立，但要求研究者报告判别证据。
\end{nlcritique}

\section{M4 校园社交：复读机与共同在场}\label{sec:07:m4}

\subsection{机制表述}

第一版教材原文：“集体居住让奶蛙拥有‘线上’‘面对面’
双通道传播途径。同一个梗在同一屋檐下被反复复诵，
趣味衰减极慢。”重写：

\begin{table}[htbp]\centering
  \caption{机制 M4（校园社交）的四列展开（\conf{M} 级）}
  \label{tab:07:m4}
  \small
  \begin{tabularx}{\textwidth}{@{}lX@{}}
    \toprule
    压力源 & 社交需求：群属感的建立与维持需要
      低成本的共同话题 \\
    补偿行为 & 群聊复读、线下引用（口头“安迪”）、
      表情包互发、立牌合影打卡 \\
    可观察指标 & 同一梗在单一群聊中的复读次数分布；
      “线上引用$\to$线下口头使用”的转化率；
      同好会纳新期的话题集中度 \\
    竞争性解释 & (a) 算法同质化：同一校园用户群本就
      被推送相似内容，复读是算法而非社交的产物；
      (b) 人口学同质：同龄同校用户的偏好本来就相近 \\
    \bottomrule
  \end{tabularx}
\end{table}

\subsection{学理对应与假设}

M4 的学理骨架是普罗文的社交笑理论
（\ref{sec:07:theory} 小节）：笑首先是一种
社交纽带行为，其发生频率取决于社交在场密度
\cite{provine2000laughter}。校园环境（集体宿舍、
固定班级、同好会组织）恰好是高社交在场密度的空间，
“线上梗 + 线下复读”的双通道因此成为可能。
第十章将展开传播网络的结构分析（弱连接与强连接
\cite{granovetter1973weak}、创新扩散的采用者类别
\cite{rogers2003diffusion}）；本章只提出
与 M4 直接相关的行为假设。

\begin{hypothesis}[H7.4 校园社交假设]\label{hyp:07:m4}
  同一奶蛙素材的“趣味衰减速度”与使用者所在群体的
  线下共同在场密度负相关：
  集体居住（宿舍）用户对该素材的持续使用时长
  显著长于非集体居住用户，且复读行为可预测
  群体凝聚力的自评指标。
\end{hypothesis}

\keyw{证伪条件}：集体与非集体居住用户的素材衰减曲线
无显著差异；或复读行为与凝聚力指标不相关。
第五章 SIR 模型的 $\gamma$（退出传播率）在 M4 框架下
获得一个机制解读：共同在场密度通过“复读强化”
降低 $\gamma$——这是把 M 级模型参数与
社会学变量挂钩的一个示例，仍然是教学性的
（\conf{M} 级），其价值在于指明该测量哪个量。

\section{M5 逆反认证：官方压制作为背书}\label{sec:07:reverse}

\subsection{机制表述}

第一版教材原文：“2026 年 3 月，奶龙官方关评、拉黑，
反而被网友读作‘官方破防’。权威的否认，
成了 meme 含金量的最高背书，热度此后一发不可收拾。”
重写：

\begin{table}[htbp]\centering
  \caption{机制 M5（逆反认证）的四列展开（\conf{M} 级）}
  \label{tab:07:m5}
  \small
  \begin{tabularx}{\textwidth}{@{}lX@{}}
    \toprule
    压力源 & 迷因社群的地位竞争：在梗快速过期的环境里，
      需要持续的“含金量”信号来区分圈内圈外 \\
    补偿行为 & 转发被官方否认的内容、以“官方破防”为题二创、
      在被拉黑后展示“阵亡”截图 \\
    可观察指标 & 官方否认事件（2026-03 关评拉黑）前后，
      奶蛙相关内容的日新增量曲线形状；
      “破防”类标题内容的占比变化 \\
    竞争性解释 & (a) 蒸汽效应（禁果效应）：单纯的好奇心驱动，
      与地位竞争无关；(b) 无关事件叠加：2026 年视频生成工具
      上线等供给侧事件才是主因（见第八章） \\
    \bottomrule
  \end{tabularx}
\end{table}

\subsection{学理对应与假设}

M5 的学理骨架是 Nissenbaum 与 Shifman 的
“迷因作为争夺中的文化资本”\cite{nissenbaum2015contested}：
在梗社群中，懂得并使用“正确”的梗是地位来源，
而\keyw{外部权威的敌意承认}是含金量最便宜的证据
——它证明该梗重要到值得被压制。巴赫金的狂欢论
提供了互补视角：权威与民间在狂欢中的等级翻转，
“官方破防”叙事本身就是一次降格表演
\cite{bakhtin1968rabelais}。

第四章已经记录了事实层面的材料
（\ref{sec:04:official}）：官方行为是
“呼吁停止丑化二创、关闭评论区、批量拉黑”，
且“直至今日你依然可以在奶龙官方视频的评论区
看见无数奶蛙图片”（\conf{C} 级）。
第五章的命题 \ref{prop:05:second}（二次爆发的条件）
则为 M5 提供了时间结构：重新激活需要
“未被耗尽的易感存量 + 足够强的重新激活项”，
官方事件恰是一个强激活项的候选。

\begin{hypothesis}[H7.5 逆反认证假设]\label{hyp:07:m5}
  2026-03 官方否认事件之后，奶蛙相关内容的日新增量
  出现\keyw{不能由供给侧事件单独解释的加速}；
  且加速期内容中，“官方破防”类框架的占比
  显著高于事件前基线。
\end{hypothesis}

\keyw{证伪条件}：控制视频生成工具上线等供给侧变量后
（第八章给出该变量的时间线），加速消失；
或加速存在但其内容中“官方破防”框架占比无变化
（说明热度上升与“否认—背书”叙事无关）。
必须强调：H7.5 的检验依赖 \conf{C} 级事件时间线
与平台数据的可得性，在两者齐备之前，
M5 的地位是\keyw{与史实相容的解释}，
不是已获支持的解释。

\begin{nlcritique}[批判性反思：逆反认证的伦理面]
  “官方压制$\to$含金量上升”若为真，它提供了一个
  令人不适的推论：\keyw{压制是迷因的养料}。
  这对权利人意味着一个两难——回应则认证，不回应则扩散；
  对研究者则意味着：M5 的验证不应变成
  “如何更好地玩逆反”的操作手册。同时必须记住
  第五章的提醒：资本框架不能取消“单纯觉得好笑”
  的朴素解释；官方否认事件在普通转发者那里
  很可能根本不构成“地位信号”，
  只是“又刷到了”。M5 预测的是
  \keyw{内容生产端}（二创作者）的框架占比变化，
  不是\keyw{消费端}的动机变化——
  把两者混为一谈是逆反认证叙事最常见的滥用。
\end{nlcritique}

\section{假设汇总与判别设计}\label{sec:07:compete}

\subsection{五个假设的一览}

\begin{table}[htbp]\centering
  \caption{五个机制假设的汇总（全部处于待检验状态；
          检验工具见第十一章）}
  \label{tab:07:hyps}
  \small
  \begin{tabularx}{\textwidth}{@{}llXX@{}}
    \toprule
    假设 & 机制 & 核心预测 & 证伪条件（简写） \\ \midrule
    H7.1 & 情绪代偿 & 压力期使用选择性上升 &
      无选择性或与自评压力无关 \\
    H7.2 & 群体稀释 & 带梗图降低删除率、
      改变回复结构 & 两指标均无差异 \\
    H7.3 & 反精致 & 敏感期劣化内容占比上升 &
      无周期性或与人设无关 \\
    H7.4 & 校园社交 & 在场密度延长素材寿命 &
      衰减曲线无居住类型差异 \\
    H7.5 & 逆反认证 & 否认事件后加速且
      “破防”框架占比上升 & 控制供给侧后加速消失 \\
    \bottomrule
  \end{tabularx}
\end{table}

\subsection{机制不是并列的：分层模型}

五个机制作用于不同层次，因此它们之间的关系
不是“五选一”，而是\keyw{分层叠加、各自可证伪}：

\begin{figure}[htbp]\centering
\begin{tikzpicture}[font=\small,
  layer/.style={rounded corners,draw=NLDeep,thick,align=center,
    minimum height=1.15cm,minimum width=4.0cm},
  lab/.style={font=\footnotesize\heiti,color=NLDeep},
  arr/.style={-{Stealth[length=2.6mm]},thick,NLDeep}]
  \node[layer,fill=NLRed!8]   (m5) at (0,0)   {M5 逆反认证\\{\footnotesize 权威关系层}};
  \node[layer,fill=NLGreen!8] (m4) at (0,-1.7){M4 校园社交\\{\footnotesize 传播层}};
  \node[layer,fill=NLViolet!8](m3) at (0,-3.4){M3 反精致\\{\footnotesize 风格层}};
  \node[layer,fill=NLAmber!12](m2) at (0,-5.1){M2 群体稀释\\{\footnotesize 个体—社交层}};
  \node[layer,fill=NLYellow!18](m1) at (0,-6.8){M1 情绪代偿\\{\footnotesize 个体—情绪层}};
  \draw[arr] (m1) -- node[right,font=\footnotesize]{情绪表达需要形式} (m2);
  \draw[arr] (m2) -- node[right,font=\footnotesize]{表达固化成风格} (m3);
  \draw[arr] (m3) -- node[right,font=\footnotesize]{风格需要传播通道} (m4);
  \draw[arr] (m4) -- node[right,font=\footnotesize]{传播引来权威反应} (m5);
  \node[lab,anchor=west] at (2.6,-3.4) {自下而上：需求逐层装配};
  \draw[arr,dashed,NLGrey] (m5.west) to[out=200,in=160] node[left,font=\footnotesize]{反馈：压制反哺下层} (m1.west);
\end{tikzpicture}
  \caption{五机制的分层装配模型（本书构造的教学框架，\conf{M} 级；
          虚线为 M5 对下层的反馈回路）。}
  \label{fig:07:layers}
\end{figure}

分层的论证意义有二。第一，\keyw{证伪是局部的}：
H7.1 被证伪不连坐 H7.3——这与第一版教材
“五重并列”的印象不同。第二，\keyw{解释有顺序}：
底层机制（M1）是发生学的起点，
顶层机制（M5）是生态的放大器；
若研究只观察到 M5（最显眼、最有戏剧性），
而跳过 M1--M4，就会把放大器误当发动机——
这是多数“官方破防”评论的结构性错误。

\subsection{竞争性解释的判别设计}

第一章保留的三种读法中，纯消遣读法与
技术—平台读法在本章换装为各机制的竞争性解释。
判别它们与功能机制的核心手法是\keyw{找不对称指标}：

\begin{enumerate}[label=\arabic*.]
  \item \textbf{周期不对称}：功能机制预测使用频率
        随压力周期波动（学期节律）；
        纯消遣与算法机制预测平稳或仅随平台活动波动。
        指标：使用频率的学期自相关。
  \item \textbf{对象不对称}：功能机制预测使用集中于
        情绪相关语境（自述、吐槽、成绩话题）；
        纯消遣预测均匀分布于一般娱乐语境。
        指标：使用语境的编码分布。
  \item \textbf{群体不对称}：功能机制预测高压力群体
        （毕业年级、考研群体）使用强度更高；
        算法机制预测使用强度仅随平台曝光量变化，
        与压力群体身份无关。
        指标：分群体的强度比较（控制曝光量）。
\end{enumerate}

三条不对称指标同时成立时，功能—符号读法
才获得超过竞争读法的证据权重；
任何一条不成立，都须相应调低对应机制的置信度。
这是“本书立场是要检验的”这句话的操作化版本
（第一章 \ref{sec:01:puzzle}）。

\section{本章小结}\label{sec:07:summary}

本章把第一版教材的五重机制重写为四列矩阵
（压力源—补偿行为—可观察指标—竞争性解释），
为每个机制补齐了学理来源（伯格森、弗洛伊德、
巴赫金、普罗文、Frijda、戈夫曼、情绪感染研究、
Hebdige、Nissenbaum--Shifman）与可检验假设
（H7.1--H7.5），并给出证伪条件。
结构上，五机制被组织为自下而上的分层装配：
情绪层（M1）$\to$社交层（M2）$\to$风格层（M3）
$\to$传播层（M4）$\to$权威关系层（M5），
M5 同时以反馈回路作用于下层。
必须重申全书纪律：本章没有任何经验数据，
五个假设全部待检验；“笑着崩溃”“我们都这样”
“官方破防”这些生动的社群表述，
在检验完成之前都只是假设的通俗名。
遗留问题有二：机制要运转，需要低成本的生产工具——
AIGC 如何把 M1--M4 的表达成本压到趋近于零，
是第八章的任务；机制要落地，需要真实的受访者——
第十一章给出 H7.1--H7.3 的问卷与 H7.4 的
行为测量设计。

\begin{nlexercises}
  \item 任选一个你熟悉的其他网络迷因，
        用表 \ref{tab:07:matrixtpl} 的四列模板为其写一个机制假设，
        并给出证伪条件。写完后自查：
        你的“可观察指标”是否能在不接触当事人隐私的前提下计数？
        若不能，请改写指标。
  \item H7.2（群体稀释）的竞争性解释之一是
        “身份认证”（用梗因为圈内人应该用梗）。
        设计一个判别实验（真实或自然实验均可），
        说明在什么结果下两个解释可以分开。
  \item 图 \ref{fig:07:layers} 的分层模型预测：
        若 M1（情绪代偿）被证伪，M3（反精致）仍可独立成立。
        请说明该预测成立与否分别对
        “奶蛙现象的可持续性”意味着什么
        （提示：情绪压力是周期性的，风格规范是长期的）。
  \item “逆反认证”若被证实，权利人将面临
        “回应即认证、不回应即扩散”的两难。
        从迷因治理的角度写 300 字分析：
        除“回应/不回应”外还存在哪些策略选项，
        各自对应本章哪个机制的变动。
\end{nlexercises}

\nlkeywords 机制矩阵、情绪代偿、群体稀释、反精致、校园社交、
逆反认证、证伪条件、竞争性解释、机械性、能量节省、降格、
社交笑、分层装配模型
